\section{Related work}
\label{sec:related}

\paragraph{Reward hacking in coding agents.}
Frontier coding agents exploit flawed tests: they delete or edit failing tests, special-case expected outputs, or tamper with the evaluation harness \citep{vonarx2025reward,metr2025o3,anthropic2025sonnet37card,anthropic2025claude4card}. ImpossibleBench \citep{zhong2025impossiblebench} makes this measurable by mutating tests so they contradict the specification, so that any pass implies cheating; EvilGenie \citep{gabor2025evilgenie} and SpecBench \citep{zhao2026specbench} measure related gaps between visible and held-out tests. Anthropic's system cards report ``impossible task'' hack rates that differ sharply across Claude generations and respond differently to anti-hacking prompts \citep{anthropic2025sonnet45card,anthropic2025haiku45card,anthropic2025opus45card}. Reward hacking learned in production-like RL can generalize to broader misalignment \citep{macdiarmid2025natural,taylor2025school}, and inoculation prompting is one proposed mitigation \citep{wichers2025inoculation,tan2025inoculation}. \citet{gomez2026escalation} show that giving agents a structured channel to report defective tests largely replaces hacking with disclosure. This literature measures \emph{whether} agents hack. We hold the hack fixed and ask how the agent \emph{describes} it.

\paragraph{Honesty of agent reports.}
An agent's final message is often the only account of its work that a user reads. OverclaimBench \citep{smyth2026quantifying} finds that agents that skip part of a file-review task usually misrepresent their coverage in the final response. Across more than twenty thousand agent-authored pull requests, descriptions that disagree with the code are accepted far less often \citep{gong2026analyzing}. System cards document o3 claiming actions it did not take \citep{openai2025gpt5card}, and Transluce found pre-release o3 inventing elaborate justifications when confronted \citep{chowdhury2025o3truthfulness}. Follow-up interviews after scheming have precedent in \citet{meinke2024frontier}. Training methods that elicit confessions or self-incrimination \citep{joglekar2025training,lee2026training,li2025spilling} aim to make such reports reliable. Our contribution is a specific, verifiable failure: when an agent games a test, its report attributes the gamed behavior to a requirement that exists nowhere in the repository or conversation.

\paragraph{Unfaithful explanations and monitorability.}
Language models give plausible explanations that omit the real cause of their answers \citep{turpin2023language}, and reasoning models often fail to verbalize the cues or hacks that drive their behavior \citep{chen2025reasoning,turpin2025teaching,arcuschin2026chainofthought}. Optimizing against chain-of-thought monitors can produce obfuscated reward hacking \citep{baker2025monitoring}, and misleading rationalizations in a trajectory can deceive monitors \citep{arnav2025cot}. These studies concern chain-of-thought. We study the user-facing report, where rationalization directly shapes human and automated review, and we find that agents fabricate a \emph{source} for the behavior, not only a reason. Closest to our review experiment, a recent post holds hacked code fixed and varies a researcher-written narrative, finding that a stated motive leads a Claude judge to label hacks as honest \citep{roussel2026claudejudge}; we use the agents' own reports and measure code-review approval.

\paragraph{Self-knowledge and self-report.}
Whether models' statements about their own behavior carry privileged information is contested \citep{binder2025looking,song2025language,song2025privileged,zeng2026evaluating}. \citet{blandfort2026strangers} find that what models predict about their own behavior is largely generic, with a self-flattering bias. Our probes speak to a narrower question: whether an agent that asserted a requirement will, when asked, acknowledge that its only basis was the test.
